\documentclass{article}
\usepackage[utf8]{inputenc}     % pdfLaTeX
\usepackage[T1]{fontenc}
\usepackage[magyar]{babel}
\usepackage{amsmath}
\usepackage{amsthm}
\usepackage[version=4]{mhchem}
\usepackage{siunitx}
\usepackage{graphicx}
% mértékegységek, hőmérséklet

\usepackage{titlesec}

\setcounter{secnumdepth}{5}
\setcounter{tocdepth}{5}

% Level 4
\titleclass{\subsubsubsection}{straight}[\subsubsection]
\newcounter{subsubsubsection}[subsubsection]
\renewcommand\thesubsubsubsection{\thesubsubsection.\arabic{subsubsubsection}}

\titleformat{\subsubsubsection}
{\normalfont\normalsize\bfseries}
{\thesubsubsubsection}{1em}{}

% Level 5
\titleclass{\subsubsubsubsection}{straight}[\subsubsubsection]
\newcounter{subsubsubsubsection}[subsubsubsection]
\renewcommand\thesubsubsubsubsection{\thesubsubsubsection.\arabic{subsubsubsubsection}}

\titleformat{\subsubsubsubsection}
{\normalfont\normalsize\bfseries}
{\thesubsubsubsubsection}{1em}{}

\titlespacing*{\subsubsubsection}{0pt}{3ex}{1ex}
\titlespacing*{\subsubsubsubsection}{0pt}{3ex}{1ex}

\title{5. Főcsoport: A Nitrogén csoport}
\author{Haydu Lőrinc, Lengyel Zoltán Péter, Rovenszky Robin}
\date{Legutóbb frissítve: 2026. Március 24.}

\newtheorem{definition}{Definíció}

\begin{document}

\maketitle
\tableofcontents
\newpage

\section{A Nitrogén csoport}
    \begin{enumerate}
        \item Nitrogén (N)
        \item Foszfor (P)
        \item Arzén (As)
        \item Antimon (Sb)
        \item Bizmut (Bi)
    \end{enumerate}
    
    5 db vegyérték elektron $\implies$ a nemesgáz elektronszerkezetének eléréséhez:
    \begin{itemize}
        \item 3 db kovalens kötés\\
        Pl.: Ammónia 
        % insert ammónia ábra here
        \item 3-szoros negatív töltésű anion (pl. $N^{3-}$ - nitrid-ion). \\
        Egyedi, mert lényegében csak a nitrogén tud ennyi iont felvenni, mivel nagyon nagy az elektronegativitása.
        \item Az is lehetséges, hogy a nitrogén 5 db kovalens kötést létesítsen. Pl.: $HNO_3$ - salétromsav
        % insert salétromsav ábra here
        \\
        A nitrogén túllépett a nemesgáz elektroszerkezetén, de ez így is energetikailag kedvező.
    \end{itemize}
    
\section{Nitrogén}
    \subsection{Atomi jellemzők}
        \begin{itemize}
            \item nagy EN (3.0)
        \end{itemize}
        
    \subsection{Molekulaszerkezet}
        %insert Nitrogen atom here
        Ez egy erős kovalens kötés $\rightarrow$ közönséges körülmények között nem reakcióképes
        $\rightarrow$ apoláris molekula $\rightarrow$ diszperziós kölcsönhatás $\rightarrow$ molekularácsos szilárd halmazállapotban
    \subsection{Fizikai tulajdonságok}
        \begin{itemize}
            \item Alacsony olvadáspont és forráspont (\SI{-210}{\degreeCelsius} és \SI{-180}{\degreeCelsius})
            \item Vízben rosszul oldódik
            \item Keszonbetegség (búvárbetegség) $\rightarrow$ gázok oldhatósága nagyobb nyomás alatt jobb!
        \end{itemize}
    
    \subsection{Kémiai tulajdonságok}
        \begin{itemize}
            \item $N_2 + O_2 \xrightarrow{\SI{3000}{\degreeCelsius}} 2NO$
            \item \(\ce{N2 + 3 H2 <=>[\text{Fe-katalizátor}] 2 NH3}\)
        \end{itemize}
    
    \subsection{Biológiai kitekintő: Nitrogén körforgása a környezetben}
        \begin{itemize}
            \item A nitrogén hármas kötése miatt majdnem semmi sem veszi fel, viszont nitrifikáló baktériumok fel tudják venni, ezek pedig szimbiózisban élnek bizonyos növényekkel (pl. bab, borsó, ill. egyéb hüvelyesek).
            \item Másik mód, amivel a körforgásba kerül a $\SI{3000}{\degreeCelsius}$-n bekövetkező reakció (villámlás).
        \end{itemize}
    
    \subsection{Szervetlen vegyipar legfontosabb folyamatai}
        \begin{itemize}
            \item Hidrogén-klorid gyártás
            \item Kénsav gyártás
            \item Ammónia gyártás
        \end{itemize}
    
    
    \subsection{Felhasználás}
        \begin{enumerate}
            \item Ammónia gyártás
            \item Védőgáz (csomagolás, hegesztés)
            \item Hűtés (folyékony nitrogénnel)
        \end{enumerate}
        
    \subsection{Előállítás} 
        A levegő cseppfolyósításával és frakcionális desztillációjával elkülönítik.
        
        \begin{definition}[Frakcionált desztilláció.]
        A forráspontkülönbségen alapuló elválasztási művelet.
        \end{definition}
    
    \subsection{A vegyületei}
        \subsubsection{Oxidjai}
            \subsubsubsection{Dinitrogén-monoxid; $N_2O$}
                    Színtelen, édeskés szagú gáz 
                    \begin{itemize}
                        \item Felhasználása: 
                        \begin{itemize}
                            \item Tejszínhab készítése (szifon)
                            \item Érzéstelenítés orvostudományban (szülési fájdalmak csökkentése)
                            \item Méhészet (elbódítja a méheket, mielőtt beavatkozik)
                            \item Autóversenyek (nitro)
                        \end{itemize}
                        \item Hétköznapi neve: kéjgáz (nevetőgáz)
                    \end{itemize}
            
            \subsubsubsection{Nitrogén-monoxid; $NO$}
                    \[
                    N_2 + O_2 \xrightarrow{\SI{3000}{\degreeCelsius}} 2NO
                    \]
                    Tovább alakul:
                    \[
                    2NO + O_2 \rightarrow 2NO_2
                    \]
                    Színtelen, szúrós szagú, mérgező gáz, amely az érfalak rugalmasságát szabályozza az emberi szervezetben, $\implies$ a vérnyomásra hat (nitroglicerin tabletta)
                    % insert nitroglicerin ábra here
            
            \subsubsubsection{Nitrogén-dioxid; $NO_2$}
                    % insert nitrogén-dioxid ábra here
                    Vörösbarna, szúrós szagú gáz, erősen mérgező. A szennyezett városi levegő egyik mérgező összetevője.
                
        \subsubsection{Oxosavai}
            \subsubsubsection{Salétromossav}
                    \[
                    2NO_2 + H_2O \rightarrow \underset{\text{salétromossav}}{HNO_2} + \underset{\text{salétromsav}}{HNO_3}
                    \]
                    Gyenge, bomlékony, egyértékű sav. Csak vizes oldatban létezik. Sói: nitritek, pl.: $KNO_2$ $\leftarrow$ kálium-nitrit $\leftarrow$ pác\textbf{só} $\leftarrow$ tartósítószer\\
            
            \subsubsubsection{Salétromsav}
                    Csak vizes oldatban létezik; erős, egyértékű sav. A koncentrált (tömény) salétromsav-oldat nagyjából 65 w\%-os.
            
                    \begin{definition}[Erős sav]
                        A $H^+$-ion leadása híg vizes oldatban 100\%-os.
                    \end{definition}
                    
                    \begin{definition}[Egyértékű sav]
                        Egy sav molekula egy db $H^+$-iont képes leadni.
                    \end{definition}
            
                    A salétromsav használható fehérjék kimutatására:\\
                    xantoprotein-próba: fehérje + cc. salétromsav
                    $\rightarrow$ dohánysárga szín\\
                    Salétromsav régies neve: választó víz, mert a $HNO_3$ oldja az ezüstöt, de nem oldja az aranyat. Ékszerészek használták.
            
                    \subsubsubsubsection{Salétromsav reakciója fémekkel}
            
                        \hrule
                        \vspace{8pt}
                    
                        \begin{enumerate}
                            \item Híg salétromsav ($\approx5$w\%-os) és ezüst
                        % Hány w%-ig?
                        
                            \setlength\parindent{110pt}
                            Nincs reakció
                            
                            \item Közepes töményénységű ($\approx$30w\%-os) salétromsav és ezüst
                            \[
                                3Ag + 4HNO_3 \rightarrow 3AgNO_3 + \underset{\text{színtelen gáz}}{NO} + 2H_2O
                            \]
                            \item Tömény (65w\%-os) salétromsav és ezüst
                            % >65w%?
                            \[
                            Ag + 2HNO_3 \rightarrow AgNO_3 + \underset{\text{vörösbarna gáz}}{NO_2} + H_2O
                            \]
        
                            \hrule
                        
                            \item Híg salétromsav ($\approx$5w\%-os) és vas
                            % Mennyi w%-ig?
                            \[
                                Fe + 2HNO_3 \rightarrow Fe(NO_3)_2 + H_2
                            \]
        
                            % közepes töménységűvel mizu?
                            
                            \item Tömény (65w\%-os) salétromsav és vas
                            % >65w%? Honnantól?
                            
                            \setlength\parindent{110pt}
                            Passziválódik!
        
                            \vspace{10pt}
                            \hrule
                            
                        \end{enumerate}
                        
                        \begin{definition}[Passziválódás]
                            Ugyan a reakció beindul, viszont egy atomnyi vastagságú "védőréteg" alakul ki, ezért a reakció megáll. A védőréteg felépítése bonyolult. A védőréteg nitrátból és oxidból áll.
                        \end{definition}
            
                    \subsubsubsubsection{A salétromsav fontos vegyületei}
                        \begin{enumerate}
                            \item $NaNO_3$; nátrium-nitrát, chilei salétrom
                            \item $KNO_3$; kálium-nitrét, kálisalétrom $\leftarrow$ fekete lőpor
                            \item $NH_4NO_3$; ammónium-nitrát $\leftarrow$ műtrágya (pétisó - $NH_4NO_3$ és mészkőpor keveréke). Nyílt lángban robbanékony.
                        \end{enumerate}

\section{Kísérletek}
    \subsection{Kísérlet 1.}
        % insert kísérlet 1 picture here
        \textbf{Leírás.} Híg salétromsav (5 w$\%$-os) + Mg-reszelék\\
        \textbf{Tapasztalat.} Pezsgés, enyhe hőfejlődés, színtelen, szagtalan gáz\\
        \textbf{Magyarázat:} 
        \[
            Mg + 2HNO_3 \rightarrow Mg (NO_3)_2 + H_2
        \]
        
    \subsection{Kísérlet 2: Tömény salétromsav és 5 Forintos érme reakciója.}
        % insert kísérlet 2 picture here
        \textbf{Leírás.} Tömény salétromsavba 5 Forintos érmét rakunk.\\
        \textbf{Tapasztalat.} Vörösbarna, jellegzetes szagú gáz, zöldeskék oldat. A főzőpohár fala forró lett, az érme jóval kisebb lett.\\
        \textbf{Magyarázat.} 
        \[
            Cu + 4HNO_3 \rightarrow \underset{\text{réz-nitrát}}{Cu(NO_3)_2} + 2NO_2 + 2H_2O
        \]
        
\end{document}